% ============================================================
% sec/01introduccion.tex
%
%Qu´e es el tema emergente, por qu´e es relevante hoy, qu´e problema de seguridad plantea
% ============================================================
\section{Introducción}
\label{sec:introduccion}

El reconocimiento facial ha pasado, en poco más de una d\'ecada, de ser
una tecnolog\'ia experimental a un componente cotidiano de control de
acceso, verificaci\'on de identidad y vigilancia p\'ublica. Aeropuertos,
bancos, tel\'efonos m\'oviles y cuerpos policiales lo han adoptado a
gran escala, impulsados por avances en aprendizaje profundo que han
llevado la exactitud de los mejores algoritmos a niveles cercanos al
100\,\% en condiciones controladas \cite{grother_nistir8280_2019}. Sin
embargo, ese mismo despliegue masivo ha expuesto un problema de
seguridad de la informaci\'on que va m\'as all\'a del rendimiento
t\'ecnico: la posibilidad de que un sistema biom\'etrico identifique
err\'oneamente a una persona, o de que sea usado para vigilarla sin su
conocimiento ni consentimiento, compromete simult\'aneamente la
integridad de la identidad atribuida a un individuo y su derecho a la
privacidad.

El problema es emergente en un doble sentido. T\'ecnicamente, porque
la evidencia acumulada muestra que la precisi\'on de estos sistemas no
es uniforme: var\'ia de forma significativa seg\'un g\'enero, edad y
origen \'etnico o geogr\'afico
\cite{grother_nistir8280_2019,kotwal_fairness_2026}, y puede ser
deliberadamente degradada mediante ataques adversarios dise\~nados
para enga\~nar al modelo \cite{vakhshiteh_adversarial_2021}.
Normativamente, porque su adopci\'on por parte de fuerzas policiales
ha superado la capacidad de los marcos regulatorios existentes para
delimitar su uso, generando pr\'acticas de vigilancia con supervisi\'on
judicial d\'ebil o inexistente
\cite{fussey_assisted_2020,almeida_ethics_2022}.

Este trabajo aborda el tema desde la tr\'iada CID: la
\emph{confidencialidad} de los datos biom\'etricos recolectados, la
\emph{integridad} de las decisiones de identificaci\'on que el sistema
entrega a un operador humano, y la \emph{disponibilidad} entendida
como la capacidad del sistema de servir por igual a toda la poblaci\'on
sobre la que se aplica. A partir de la revisi\'on de cinco fuentes que
combinan un informe t\'ecnico de referencia (NIST), estudios emp\'iricos
de despliegue policial, an\'alisis regulatorio comparado y encuestas
acad\'emicas sobre sesgo y ataques adversarios, se construye un
panorama del estado del arte que conecta directamente con el debate de
privacidad discutido en el curso y anticipa los contenidos de
biometr\'ia que se abordar\'an m\'as adelante.
